\documentclass[12pt,a4paper]{article}
\usepackage[UTF8]{ctex}
\usepackage{amsmath,amssymb,amsthm}
\usepackage{geometry}
\usepackage{graphicx}
\usepackage{hyperref}
\usepackage{bm}

\geometry{margin=2.5cm}

\title{质量矩阵公式$M(\theta) = \sum_{i=1}^{n} J_i^b(\theta)^T G_i J_i^b(\theta)$的详细推导}
\author{Modern Robotics}
\date{}

\begin{document}

\maketitle

\tableofcontents
\newpage

\section{引言}

在机器人动力学中，质量矩阵$M(\theta)$是一个关键概念，它将关节加速度$\ddot{\theta}$与关节力矩$\tau$联系起来。质量矩阵的一个重要表达式是：

\begin{equation}
M(\theta) = \sum_{i=1}^{n} J_i^b(\theta)^T G_i J_i^b(\theta)
\label{eq:mass_matrix}
\end{equation}

本文将详细解释为什么质量矩阵可以这样表示，以及这个公式的物理和数学意义。

\section{推导的起点：动能表达式}

\subsection{总动能为各连杆动能之和}

机器人的总动能$K$等于所有连杆的动能之和：

\begin{equation}
K = \sum_{i=1}^{n} K_i
\label{eq:total_kinetic_energy}
\end{equation}

其中$K_i$是连杆$i$的动能。

\subsection{单个连杆的动能}

对于连杆$i$，其动能可以表示为：

\begin{equation}
K_i = \frac{1}{2} V_i^T G_i V_i
\label{eq:link_kinetic_energy}
\end{equation}

其中：
\begin{itemize}
\item $V_i \in \mathbb{R}^6$：连杆$i$的twist（在坐标系$\{i\}$中表示）
\item $G_i \in \mathbb{R}^{6 \times 6}$：连杆$i$的空间惯性矩阵（在坐标系$\{i\}$中表示）
\end{itemize}

\textbf{为什么是这种形式？}

空间惯性矩阵$G_i$的定义使得：
\begin{equation}
G_i = \begin{bmatrix} I_i & 0 \\ 0 & m_i I \end{bmatrix}
\end{equation}

其中$I_i$是连杆$i$的惯性张量，$m_i$是质量。因此：
\begin{align}
V_i^T G_i V_i &= \begin{bmatrix} \omega_i^T & v_i^T \end{bmatrix} \begin{bmatrix} I_i & 0 \\ 0 & m_i I \end{bmatrix} \begin{bmatrix} \omega_i \\ v_i \end{bmatrix} \\
&= \omega_i^T I_i \omega_i + m_i v_i^T v_i
\end{align}

这正是旋转动能$\frac{1}{2}\omega_i^T I_i \omega_i$和平移动能$\frac{1}{2}m_i v_i^T v_i$之和。

\textbf{因此}：
\begin{equation}
K = \frac{1}{2}\sum_{i=1}^{n} V_i^T G_i V_i
\label{eq:kinetic_energy_sum}
\end{equation}

\section{关键步骤：将Twist表示为关节速度的函数}

\subsection{Body雅可比的定义}

设$T_{0i}(\theta_1, \ldots, \theta_i)$表示从基座坐标系$\{0\}$到连杆坐标系$\{i\}$的正向运动学变换。

\textbf{Body雅可比}$J_i^b(\theta)$定义为：从$T_{0i}^{-1}\dot{T}_{0i}$提取的雅可比矩阵，它将关节速度$\dot{\theta}$映射到连杆$i$的Body Twist：

\begin{equation}
V_i = J_i^b(\theta) \dot{\theta}
\label{eq:body_jacobian}
\end{equation}

\textbf{注意}：
\begin{itemize}
\item $J_i^b(\theta)$是一个$6 \times n$矩阵（$n$是总关节数）
\item 前$i$列对应影响连杆$i$的关节
\item 后$n-i$列全为零（因为后$n-i$个关节不影响连杆$i$）
\end{itemize}

\subsection{为什么可以这样表示？}

\textbf{物理意义}：连杆$i$的速度$V_i$是由前$i$个关节的运动产生的。每个关节$j$（$j = 1, \ldots, i$）对连杆$i$速度的贡献是$A_j^b \dot{\theta}_j$，其中$A_j^b$是关节$j$的螺旋轴在坐标系$\{i\}$中的表示。

因此：
\begin{equation}
V_i = \sum_{j=1}^{i} A_j^b \dot{\theta}_j = J_i^b(\theta) \dot{\theta}
\end{equation}

其中$J_i^b(\theta)$的第$j$列是$A_j^b$（$j \leq i$），否则为零。

\section{推导质量矩阵公式}

\subsection{步骤1：代入Body雅可比}

将方程(\ref{eq:body_jacobian})代入方程(\ref{eq:kinetic_energy_sum})：

\begin{align}
K &= \frac{1}{2}\sum_{i=1}^{n} V_i^T G_i V_i \\
&= \frac{1}{2}\sum_{i=1}^{n} (J_i^b(\theta) \dot{\theta})^T G_i (J_i^b(\theta) \dot{\theta}) \\
&= \frac{1}{2}\sum_{i=1}^{n} \dot{\theta}^T J_i^b(\theta)^T G_i J_i^b(\theta) \dot{\theta}
\label{eq:kinetic_energy_jacobian}
\end{align}

\textbf{矩阵转置的性质}：$(AB)^T = B^T A^T$，因此$(J_i^b \dot{\theta})^T = \dot{\theta}^T (J_i^b)^T$。

\subsection{步骤2：提取$\dot{\theta}$}

由于$\dot{\theta}$是公共因子，可以提取出来：

\begin{align}
K &= \frac{1}{2}\sum_{i=1}^{n} \dot{\theta}^T J_i^b(\theta)^T G_i J_i^b(\theta) \dot{\theta} \\
&= \frac{1}{2}\dot{\theta}^T \left(\sum_{i=1}^{n} J_i^b(\theta)^T G_i J_i^b(\theta)\right) \dot{\theta}
\label{eq:kinetic_energy_extracted}
\end{align}

\textbf{关键观察}：括号内的项$\sum_{i=1}^{n} J_i^b(\theta)^T G_i J_i^b(\theta)$是一个$n \times n$矩阵。

\subsection{步骤3：动能的一般形式}

在拉格朗日动力学中，动能$K$必须是关于广义速度$\dot{\theta}$的\textbf{二次型}：

\begin{equation}
K = \frac{1}{2}\dot{\theta}^T M(\theta) \dot{\theta}
\label{eq:kinetic_energy_general}
\end{equation}

其中$M(\theta)$是\textbf{质量矩阵}（也称为广义质量矩阵或惯性矩阵）。

\textbf{为什么必须是二次型？}

\begin{enumerate}
\item \textbf{物理原因}：动能是速度的平方的函数（$K = \frac{1}{2}mv^2$），因此必须是速度的二次型。

\item \textbf{数学原因}：在拉格朗日方程中，动能对广义速度的偏导数是：
\begin{equation}
\frac{\partial K}{\partial \dot{\theta}} = M(\theta) \dot{\theta}
\end{equation}

这要求$K$是$\dot{\theta}$的二次型。

\item \textbf{能量守恒}：动能必须是正定的二次型，这保证了系统的能量是正的。
\end{enumerate}

\subsection{步骤4：识别质量矩阵}

比较方程(\ref{eq:kinetic_energy_extracted})和(\ref{eq:kinetic_energy_general})：

\begin{align}
K &= \frac{1}{2}\dot{\theta}^T \left(\sum_{i=1}^{n} J_i^b(\theta)^T G_i J_i^b(\theta)\right) \dot{\theta} \\
&= \frac{1}{2}\dot{\theta}^T M(\theta) \dot{\theta}
\end{align}

\textbf{因此}，质量矩阵必须是：

\begin{equation}
M(\theta) = \sum_{i=1}^{n} J_i^b(\theta)^T G_i J_i^b(\theta)
\label{eq:mass_matrix_final}
\end{equation}

\textbf{证毕！}

\section{公式的物理意义}

\subsection{每一项的贡献}

质量矩阵$M(\theta)$的每一项$J_i^b(\theta)^T G_i J_i^b(\theta)$表示：

\begin{itemize}
\item \textbf{连杆$i$对总惯性的贡献}
\item 通过Body雅可比$J_i^b$将连杆$i$的惯性$G_i$从twist空间映射到关节空间
\item 所有连杆的贡献相加得到总的质量矩阵
\end{itemize}

\subsection{矩阵乘法的意义}

\textbf{$J_i^b(\theta)^T G_i J_i^b(\theta)$的含义}：

\begin{enumerate}
\item \textbf{$J_i^b(\theta) \dot{\theta}$}：将关节速度映射到连杆$i$的twist
\item \textbf{$G_i (J_i^b(\theta) \dot{\theta})$}：计算连杆$i$的动量（twist乘以惯性矩阵）
\item \textbf{$(J_i^b(\theta) \dot{\theta})^T G_i (J_i^b(\theta) \dot{\theta})$}：计算连杆$i$的动能（twist的转置乘以动量）
\item \textbf{提取$\dot{\theta}$}：得到关于$\dot{\theta}$的二次型，系数矩阵就是$J_i^b(\theta)^T G_i J_i^b(\theta)$
\end{enumerate}

\subsection{质量矩阵的性质}

\begin{enumerate}
\item \textbf{对称性}：$M(\theta) = M(\theta)^T$

\textbf{证明}：
\begin{align}
M(\theta)^T &= \left(\sum_{i=1}^{n} J_i^b(\theta)^T G_i J_i^b(\theta)\right)^T \\
&= \sum_{i=1}^{n} (J_i^b(\theta)^T G_i J_i^b(\theta))^T \\
&= \sum_{i=1}^{n} J_i^b(\theta)^T G_i^T (J_i^b(\theta)^T)^T \\
&= \sum_{i=1}^{n} J_i^b(\theta)^T G_i J_i^b(\theta) \quad \text{（因为$G_i$对称）} \\
&= M(\theta)
\end{align}

\item \textbf{正定性}：$M(\theta)$是正定矩阵

\textbf{证明}：对于任意非零$\dot{\theta}$：
\begin{align}
\dot{\theta}^T M(\theta) \dot{\theta} &= \dot{\theta}^T \left(\sum_{i=1}^{n} J_i^b(\theta)^T G_i J_i^b(\theta)\right) \dot{\theta} \\
&= \sum_{i=1}^{n} (J_i^b(\theta) \dot{\theta})^T G_i (J_i^b(\theta) \dot{\theta}) \\
&= 2\sum_{i=1}^{n} K_i > 0
\end{align}

因为每个连杆的动能$K_i > 0$（除非所有速度为零），所以$M(\theta)$是正定的。

\item \textbf{配置依赖性}：$M(\theta)$依赖于关节配置$\theta$，因为Body雅可比$J_i^b(\theta)$依赖于$\theta$。
\end{enumerate}

\section{具体例子}

\subsection{例子1：单连杆机器人}

对于单连杆机器人（$n = 1$）：

\begin{equation}
M(\theta) = J_1^b(\theta)^T G_1 J_1^b(\theta)
\end{equation}

如果关节是旋转关节，$J_1^b$是$6 \times 1$矩阵，$M(\theta)$是标量（$1 \times 1$矩阵），即等效转动惯量。

\subsection{例子2：2R平面机器人}

对于2R平面机器人（$n = 2$）：

\begin{align}
M(\theta) &= J_1^b(\theta)^T G_1 J_1^b(\theta) + J_2^b(\theta)^T G_2 J_2^b(\theta) \\
&= \begin{bmatrix} M_{11} & M_{12} \\ M_{21} & M_{22} \end{bmatrix}
\end{align}

其中：
\begin{itemize}
\item $M_{11}$：关节1的等效惯量（包括连杆1和连杆2的影响）
\item $M_{22}$：关节2的等效惯量（主要是连杆2的影响）
\item $M_{12} = M_{21}$：关节1和关节2之间的耦合惯量
\end{itemize}

\section{与拉格朗日方程的关系}

\subsection{拉格朗日方程}

在拉格朗日动力学中，运动方程为：

\begin{equation}
\tau = \frac{d}{dt}\frac{\partial K}{\partial \dot{\theta}} - \frac{\partial K}{\partial \theta}
\end{equation}

\subsection{计算$\frac{\partial K}{\partial \dot{\theta}}$}

使用$K = \frac{1}{2}\dot{\theta}^T M(\theta) \dot{\theta}$：

\begin{align}
\frac{\partial K}{\partial \dot{\theta}} &= \frac{\partial}{\partial \dot{\theta}} \left(\frac{1}{2}\dot{\theta}^T M(\theta) \dot{\theta}\right) \\
&= M(\theta) \dot{\theta}
\end{align}

\textbf{因此}：
\begin{equation}
\frac{d}{dt}\frac{\partial K}{\partial \dot{\theta}} = M(\theta) \ddot{\theta} + \dot{M}(\theta) \dot{\theta}
\end{equation}

这给出了动力学方程中的$M(\theta)\ddot{\theta}$项和$\dot{M}(\theta)\dot{\theta}$项（后者是科里奥利力的一部分）。

\section{总结}

\subsection{推导的关键步骤}

\begin{enumerate}
\item \textbf{总动能}：$K = \frac{1}{2}\sum_{i=1}^{n} V_i^T G_i V_i$

\item \textbf{Body雅可比}：$V_i = J_i^b(\theta) \dot{\theta}$

\item \textbf{代入}：$K = \frac{1}{2}\dot{\theta}^T \left(\sum_{i=1}^{n} J_i^b(\theta)^T G_i J_i^b(\theta)\right) \dot{\theta}$

\item \textbf{识别质量矩阵}：由于$K = \frac{1}{2}\dot{\theta}^T M(\theta) \dot{\theta}$，所以：
\begin{equation}
M(\theta) = \sum_{i=1}^{n} J_i^b(\theta)^T G_i J_i^b(\theta)
\end{equation}
\end{enumerate}

\subsection{为什么这个公式是正确的？}

\begin{itemize}
\item \textbf{数学上}：动能必须是关于$\dot{\theta}$的二次型，系数矩阵就是质量矩阵
\item \textbf{物理上}：质量矩阵将关节空间的加速度与力矩联系起来，而每个连杆通过其Body雅可比贡献惯性
\item \textbf{几何上}：Body雅可比将关节速度映射到twist空间，惯性矩阵在twist空间中作用，然后通过转置的雅可比映射回关节空间
\end{itemize}

\subsection{公式的重要性}

\begin{itemize}
\item 提供了质量矩阵的\textbf{解析表达式}
\item 揭示了质量矩阵的\textbf{结构}：各连杆贡献的叠加
\item 便于\textbf{数值计算}：可以递归计算每个连杆的贡献
\item 揭示了\textbf{耦合项}的来源：不同连杆的雅可比之间的相互作用
\end{itemize}

\section{参考文献}

\begin{itemize}
\item Modern Robotics: Mechanics, Planning, and Control
\item 第8章：开链动力学
\item 拉格朗日动力学理论
\end{itemize}

\end{document}

